\section{A worked example first}\label{sec:example}

Before defining anything, we look at a small calculation that has all the features QAM is meant to manage: two frames of description,
a map between them, an assumption about the set-up, a measurement, and an exact result. The physics is elementary and fully derived
here; the reader needs only the transformation law of electric and magnetic fields under a Lorentz boost. We use units with $c=1$.

\subsection{A field configuration and its Gram register}

In a laboratory frame, which we call \emph{Lab}, consider an electric field and a magnetic field of the form
\begin{equation}\label{eq:prep}
E = 2\,e_x,\qquad B = t\,e_y + u\,e_z,\qquad t^2+u^2 = 1 .
\end{equation}
The magnetic field has unit length and lies in the $yz$-plane; the number $t\in[-1,1]$ fixes its orientation. Put the two fields
side by side as the columns of a $3\times2$ matrix $F=[E\mid B]$. Its Gram matrix and the normalised Gram matrix are
\[
G = F^{\mathsf T}F = \begin{pmatrix} |E|^2 & E\cdot B\\ E\cdot B & |B|^2\end{pmatrix},\qquad
\rho = \frac{G}{\tr G}.
\]
We call $\rho$ the \emph{Gram register}: a positive matrix with unit trace that records the relative strengths of the two fields and
the angle between them, but forgets their overall scale and their orientation in space. For every field~\eqref{eq:prep} we have
$E\cdot B=0$, $|E|^2=4$ and $|B|^2=1$, so
\begin{equation}\label{eq:rho-lab}
\rho = \diag\bigl(\tfrac45,\tfrac15\bigr)\qquad\text{for every } t\in[-1,1].
\end{equation}

\subsection{The same fields seen by a moving observer}

A second observer moves with speed $v$ along the $z$-axis. We call that observer's frame \emph{Boosted}. With
$\gamma=(1-v^2)^{-1/2}$ the fields transform as
\[
\begin{aligned}
E'_x&=\gamma(E_x-vB_y), & E'_y&=\gamma(E_y+vB_x), & E'_z&=E_z,\\
B'_x&=\gamma(B_x+vE_y), & B'_y&=\gamma(B_y-vE_x), & B'_z&=B_z .
\end{aligned}
\]
For the fields~\eqref{eq:prep} this gives
\[
E' = \gamma\,(2-vt)\,e_x,\qquad B' = \gamma\,(t-2v)\,e_y + u\,e_z .
\]
We fix $v=\tfrac14$, so $\gamma^2=\tfrac{16}{15}$. The two boosted fields are still orthogonal. With
\[
A(t) = \bigl(2-\tfrac t4\bigr)^2
\]
we find $|E'|^2=\gamma^2A(t)$ and, using $u^2=1-t^2$,
\[
|B'|^2 = \gamma^2\bigl(t-\tfrac12\bigr)^2 + 1 - t^2 = \gamma^2\Bigl(A(t)-\tfrac{45}{16}\Bigr).
\]
The common factor $\gamma^2$ cancels in the normalisation, and the Gram register in the boosted frame is
\begin{equation}\label{eq:rho-boosted}
\rho' = \diag(p,\,1-p),\qquad p(t)=\frac{A(t)}{2A(t)-\frac{45}{16}}=\frac{(8-t)^2}{2(8-t)^2-45}.
\end{equation}
The denominator is at least $2\cdot 49-45=53$, so $p$ is well defined on $[-1,1]$.

\subsection{The point of the example}

Evaluate~\eqref{eq:rho-boosted} at the two ends of the range. For $t=1$ we get $p=49/53$, and for $t=-1$ we get $p=81/117=9/13$.
Both fields had the \emph{same} Gram register~\eqref{eq:rho-lab} in the Lab frame, but they have \emph{different} Gram registers in
the Boosted frame. Hence there is no function $f$ with $\rho'=f(\rho)$: a map from Lab registers to Boosted registers does not exist.
The boost acts on the fields; the register has forgotten the orientation that the boost needs.

This is precisely the kind of fact that is easy to lose in prose. A sentence such as ``the boost maps the Lab register to the Boosted
register'' sounds harmless and is false. In QAM, the map between the two frames is declared on the full description of the
fields, and every statement about the Boosted frame names the Lab statement and the map it came from. Section~\ref{sec:frames}
shows how; it also shows how adding the missing orientation and scale to the register restores a well-defined map.

\subsection{A measurement narrows the prediction}

Suppose we know nothing about $t$. Then all values of $p$ between the two endpoints can occur. To see that they all do, note that
$A$ decreases as $t$ grows (since $2-t/4>0$) and that
\[
\frac{dp}{dA} = -\frac{45/16}{\bigl(2A-45/16\bigr)^2} < 0,
\]
so $p$ is continuous and strictly increasing in $t$. Its range over $t\in[-1,1]$ is therefore the interval $[9/13,\,49/53]$, of width
about $0.232$.

Now suppose a measurement in the Lab frame tells us that $t=B\cdot e_y$ lies in $[-\tfrac14,\tfrac14]$. By monotonicity, the
prediction for the boosted frame is the image of the endpoints:
\begin{equation}\label{eq:interval}
p(-\tfrac14)=\frac{(33/4)^2}{2(33/4)^2-45}=\frac{121}{162},\qquad
p(\tfrac14)=\frac{(31/4)^2}{2(31/4)^2-45}=\frac{961}{1202},
\end{equation}
so $p\in[121/162,\,961/1202]$, an interval of width about $0.0526$, less than a quarter of the original one. This is a deterministic
constraint: no probability distribution over $t$ is assumed.

\subsection{What the example asks of a record}

The calculation is short, but a record of it has to keep several things apart:
\begin{itemize}[leftmargin=1.5em]
\item The preparation~\eqref{eq:prep} and the boost speed $v=\frac14$ are \emph{assumptions} about the set-up. Every conclusion
holds only under them.
\item The interval $t\in[-\frac14,\frac14]$ is a \emph{measured input}. The narrowed prediction~\eqref{eq:interval} rests on it and
must never be cited as if it followed from the theory alone.
\item Statements~\eqref{eq:rho-lab} and~\eqref{eq:rho-boosted} live in \emph{different frames}. They may be compared only through
the declared map that the boost provides, and that map acts on fields, not on registers.
\end{itemize}
These are exactly the three silent transitions of Section~\ref{sec:intro}. The rest of the article shows how QAM records them and
what it can and cannot guarantee about the record.
